\clearpage\subsection{Отчёт и защита работы №\,2}
\subsection*{Содержание отчёта}
\begin{enumerate}
\item ФИО, ИТМО ID, название и номер работы.
\item Цель, характеристика исходных данных и смысл используемых признаков.
\item Описание разбиения и предобработки: что удалено, как заполнены
пропуски, как обработаны категории и где вычислены статистики.
\item Архитектура исходной сети, размерности и расчёт числа параметров.
\item Таблица базовых прогнозов и всех выполненных конфигураций.
\item Кривые обучения, сравнение фактических и предсказанных цен,
анализ остатков и наиболее крупных ошибок.
\item Вывод о влиянии исследованных факторов и обоснование выбранной
конфигурации с указанием ограничений.
\item Пример итоговой таблицы прогнозов и описание файлов результата.
\item Список источников и параметры воспроизводимости.
\end{enumerate}

К отчёту приложите ноутбук с выполненными ячейками, исходный код изменений,
журнал опытов, истории обучения и прогнозы. Сохраняйте результаты всех
запланированных запусков, включая неудачные, с объяснением причины сбоя.
Полные CSV не требуется вставлять в PDF отчёта.

\subsection*{Критерии оценки}
\noindent\begin{tabularx}{\textwidth}{@{}p{0.28\textwidth}Y@{}}
\toprule
Критерий & Что необходимо показать \\
\midrule
Корректность данных & Разделение до обучения преобразований,
нет \file{Id} и цели во входах, согласованное кодирование. \\
Корректность модели & Подходящие вход и выход, согласованная потеря,
понимание количества параметров и процесса обучения. \\
Качество исследования & Сопоставимые опыты, сохранённые условия,
одинаковые метрики, объяснение результатов. \\
Анализ и вывод & Разбор ошибок и ограничений; вывод подтверждён
собственными наблюдениями. \\
Воспроизводимость & Данные, версии, разбиение, конфигурации и материалы
позволяют повторить опыт. \\
\bottomrule
\end{tabularx}

Баллы, пороги и правила повторной сдачи устанавливает преподаватель
по рабочей программе дисциплины. Исправный запуск кода без анализа
не демонстрирует всех результатов обучения этой работы.

\Needspace{9\baselineskip}
\subsection*{Контрольные вопросы}
\begin{enumerate}
\item Почему прогнозирование \file{SalePrice} относится к регрессии?
Чем признак отличается от класса и целевой переменной?
\item Для чего нужны обучение, валидация и \file{test.csv}?
На какой части допустим подбор гиперпараметров?
\item Почему нельзя отдельно обучать кодировщики категорий на обучающих
и новых данных? Приведите пример ошибки соответствия кодов.
\item Зачем масштабировать признаки и целевую цену? Как вернуть прогноз
к исходным единицам?
\item Что вычисляет полносвязный слой? Как найти число его параметров?
\item Почему входная размерность не меняется вместе с шириной скрытого слоя?
\item Что такое функция активации? Сравните ReLU, sigmoid и tanh.
Почему выход данной сети линейный?
\item Что такое эпоха и итерация? Сколько обновлений весов было
в вашем опыте и как на это повлиял \file{batch\_size}?
\item Что измеряют MSE, MAE и RMSE? В каких единицах они выражаются?
\item Почему при смене функции потерь нужно сохранять общие метрики сравнения?
\item Как ширина, глубина, число эпох, размер пакета и оптимизатор
повлияли на ваши результаты? Какие выводы поддерживаются данными?
\item Какие признаки переобучения вы наблюдали? Как работает ранний останов?
\item Почему лучший результат на одной валидации не гарантирует
лучшего результата на будущих домах?
\item Можно ли судить о качестве прогноза по первым строкам таблицы
\file{Id, SalePrice}, если правильная цена неизвестна?
\end{enumerate}

\begin{keybox}{Перед сдачей}
Убедитесь, что ноутбук выполняется последовательно из чистого состояния,
в отчёте есть собственные результаты всех обязательных сравнений,
а файл прогнозов содержит ровно те \file{Id} и в том же порядке,
что исходный \file{test.csv}.
\end{keybox}
